\chapter{欧几里得空间}\label{ch:b1:euclid}

在实向量空间上引入内积后，即得到几何概念——
长度、夹角、正交性与距离，以及本章的核心定理：
每个子空间都存在正交投影，
可由格拉姆–施密特方法计算，它给出最短距离。
平面上的等距变换为本章作结，
也为本学年的几何内容收尾。

全章中，$E$ 为\emph{实}向量空间。

\section{内积}

\begin{definition}\label{def:b1:euclid:def}
称 $E$ 上的\emph{内积}\index{内积}为满足下列条件的映射
$\langle\cdot,\cdot\rangle \colon E \times E \to \R$，它
是双线性的、对称的，且正定的（$\langle x, x\rangle >
0$，当 $x \neq 0$ 时成立）。配备此内积的有限维空间称为
\emph{欧几里得空间}\index{欧几里得空间}。$x$ 的\emph{范数}
为 $\norm{x} = \sqrt{\langle x, x\rangle}$，并令 $d(x, y) =
\norm{x - y}$。
\end{definition}

\begin{example}\label{ex:b1:euclid:examples}
在 $\R^n$ 上：标准内积 $\langle x, y\rangle = \sum x_i
y_i$。在 $C(\intcc{a}{b})$ 上：$\langle f, g \rangle = \int_a^b fg$
（正定性由\cref{thm:b1:integration:props}~(4)可得）。
在 $\R_n[X]$ 上：$\langle P, Q\rangle = \int_0^1 PQ$，或 $\sum_{i}
P(x_i)Q(x_i)$，对 $n+1$ 个互异点求和。
\end{example}

\begin{example}[两个多项式之间的夹角]
\label{ex:b1:euclid:polyangle}
一旦选定内积，\emph{任意}两个非零向量
之间都有夹角，由 $\cos\theta = \frac{\langle x,
y\rangle}{\norm x\,\norm y}$（由柯西–施瓦茨不等式，
该值为合法的余弦值）。对 $X$ 与 $X^2$，取 $\int_0^1$ 上的内积：
\[
\langle X, X^2\rangle = \frac14, \qquad
\norm X = \frac1{\sqrt3}, \qquad \norm{X^2} = \frac1{\sqrt5},
\qquad
\cos\theta = \frac{1/4}{1/\sqrt{15}} = \frac{\sqrt{15}}{4}
\approx 0.968 :
\]
夹角约为 $14.5$ 度 --- 在 $\intcc{0}{1}$ 上，
$x$ 与 $x^2$ 的图像在平方平均意义下“近乎平行”，这正是
为何去掉这一公共方向（见下文格拉姆–施密特方法）
后只剩下较小的修正
项 $X^2 - X + \frac16$。
\end{example}

\begin{theorem}[柯西–施瓦茨不等式；范数的性质]\label{thm:b1:euclid:cs}
对所有 $x, y \in E$：\index{柯西–施瓦茨不等式}
\[
\abs{\langle x, y\rangle} \leq \norm x\, \norm y ,
\]
当且仅当 $x, y$ 成比例时取等号。由此，$\norm\cdot$
满足三角不等式 $\norm{x + y} \leq \norm x + \norm
y$（并且 $\norm{\lambda x} = \abs\lambda \norm x$，$\norm x = 0 \iff
x = 0$）。此外：
\[
\norm{x+y}^2 = \norm x^2 + 2\langle x, y\rangle + \norm y^2,
\qquad
\langle x, y \rangle = \tfrac14\bigl(\norm{x+y}^2 -
\norm{x-y}^2\bigr) .
\]
\end{theorem}

\begin{proof}
若 $y = 0$，结论显然。否则二次式 $t
\mapsto \norm{x + ty}^2 = \norm x^2 + 2t\langle x, y\rangle + t^2
\norm y^2$ 满足 $\geq 0$（对任意 $t$），故其判别式 $\leq 0$，
这即柯西–施瓦茨不等式；等号成立意味着二重根 $t_0$，亦即
$x + t_0 y = 0$（由正定性），即成比例。三角
不等式：展开，
\[
\norm{x + y}^2 = \norm x^2 + 2\langle x, y\rangle + \norm y^2
\leq \norm x^2 + 2\norm x\,\norm y + \norm y^2
= \bigl(\norm x + \norm y\bigr)^2 ,
\]
其中中间一步用到柯西–施瓦茨不等式；等号成立迫使 $\langle
x, y\rangle = \norm x\norm y$，即\emph{正}取等情形，
亦即比值非负的成比例 ---
几何上，三角形仅当两向量
指向同一方向时才退化。最后两个
等式是直接展开的（第二式即
\emph{极化恒等式}，它由范数恢复内积）。
\end{proof}

\section{正交性}

\begin{definition}\label{def:b1:euclid:orthogonal}
$x \perp y$ 是指 $\langle x, y \rangle = 0$。称一个向量族为
\emph{正交的}，当其中向量两两正交时；
若其中每个向量还满足
范数为 $1$，则称其为\emph{标准正交族}\index{标准正交族}。子空间 $F$ 的\emph{正交补}为
\[
F^{\perp} = \{x \in E : \forall y \in F,\ \langle x, y\rangle =
0\},
\]
是 $E$ 的子空间。\index{正交补}
\end{definition}

\begin{proposition}\label{prop:b1:euclid:orthfree}
（毕达哥拉斯\index{毕达哥拉斯定理}）若 $x \perp y$，则
$\norm{x+y}^2 = \norm x^2 + \norm y^2$。非零正交向量族
是线性无关的。在标准正交基 $(e_1, \dots,
e_n)$ 中，坐标与内积分别为
\[
x = \sum_{i} \langle x, e_i\rangle\, e_i,
\qquad
\langle x, y \rangle = \sum_i \langle x, e_i\rangle \langle y,
e_i\rangle,
\qquad
\norm x^2 = \sum_i \langle x, e_i\rangle^2 .
\]
\end{proposition}

\begin{proof}
勾股定理：展开。线性无关：对零线性组合取 $\langle\,\cdot\,, x_j\rangle$，
得 $\lambda_j \norm{x_j}^2 = 0$。坐标：
设 $x = \sum \lambda_i e_i$，与 $e_j$ 作内积：
$\lambda_j = \langle x, e_j\rangle$；两式由
双线性性即得。
\end{proof}

\begin{example}[标准正交坐标，附帕塞瓦尔检验]
\label{ex:b1:euclid:parsevalcheck}
将 $x = (1, 2, 3)$ 在
\cref{exo:b1:euclid:3}的标准正交基下展开：
\[
e_1 = \tfrac{1}{\sqrt2}(1,1,0), \quad
e_2 = \tfrac{1}{\sqrt6}(1,-1,2), \quad
e_3 = \tfrac{1}{\sqrt3}(-1,1,1).
\]
无需解方程组 --- 只需三个内积：
\[
\langle x, e_1\rangle = \frac{3}{\sqrt2},
\qquad
\langle x, e_2\rangle = \frac{1 - 2 + 6}{\sqrt6} =
\frac{5}{\sqrt6},
\qquad
\langle x, e_3\rangle = \frac{-1 + 2 + 3}{\sqrt3} =
\frac{4}{\sqrt3}.
\]
用命题中的范数公式检验：
\[
\frac{9}{2} + \frac{25}{6} + \frac{16}{3}
= \frac{27 + 25 + 32}{6} = 14 = \norm{x}^2 = 1 + 4 + 9 .
\]
这一坐标平方和检验（即有限形式的帕塞瓦尔
恒等式）只需几秒，几乎必能发现符号与归一化
错误 --- 每当算出标准正交展开，都应养成这一习惯；
其无穷维版本，即
针对傅里叶系数的情形（
见\cref{ex:b1:euclid:trigortho}），
是本科三年级分册中的一条定理。
\end{example}

\begin{theorem}[格拉姆–施密特]\label{thm:b1:euclid:gramschmidt}
每个欧几里得空间都有标准正交基。具体地，对任一
基 $(v_1, \dots, v_n)$，递推式\index{格拉姆–施密特过程}
\[
w_k = v_k - \sum_{i=1}^{k-1} \langle v_k, e_i\rangle\, e_i ,
\qquad
e_k = \frac{w_k}{\norm{w_k}}
\]
给出标准正交基 $(e_1, \dots, e_n)$，使得
$\operatorname{Vect}(e_1, \dots, e_k) = \operatorname{Vect}(v_1,
\dots, v_k)$ 对每个 $k$ 成立。
\end{theorem}

\begin{proof}
对 $k$ 归纳。设 $(e_1, \dots, e_{k-1})$ 为标准正交组，
且张成 $\operatorname{Vect}(v_1, \dots, v_{k-1})$。向量
$w_k$ 由构造与每个 $e_j$（$j < k$）正交，
($\langle w_k, e_j \rangle = \langle v_k, e_j\rangle - \langle v_k,
e_j\rangle$），且 $w_k \neq 0$，因为 $v_k \notin
\operatorname{Vect}(v_1, \dots, v_{k-1})$。归一化保持
正交性；张成关系成立，因为 $e_k$ 是
$v_k$ 与前面的 $e_i$ 的可逆线性组合。
\end{proof}

\begin{example}[多项式上的格拉姆–施密特过程，完整演示]
\label{ex:b1:euclid:legendre}
对 $(1, X, X^2)$ 在 $\R_2[X]$ 中作标准正交化，内积为 $\langle P,
Q\rangle = \int_0^1 PQ$。\emph{第一步}：$\norm{1}^2 = 1$，故
$e_1 = 1$。\emph{第二步}：$w_2 = X - \langle X, 1\rangle\,1 = X -
\frac12$，且 $\norm{w_2}^2 = \int_0^1\bigl(x - \frac12\bigr)^2
\dd x = \frac1{12}$：$e_2 = \sqrt{12}\,\bigl(X - \frac12\bigr)$。
\emph{第三步}：$\langle X^2, e_1\rangle = \frac13$，又
\[
\langle X^2, e_2\rangle
= \sqrt{12}\int_0^1 x^2\Bigl(x - \frac12\Bigr)\dd x
= \frac{\sqrt{12}}{12},
\qquad\text{故}\qquad
w_3 = X^2 - \frac13 - \Bigl(X - \frac12\Bigr)
= X^2 - X + \frac16 .
\]
其范数已在\cref{exo:b1:euclid:9}中算出：$\norm{w_3}^2 =
\frac1{180}$，由此 $e_3 = \sqrt{180}\,\bigl(X^2 - X +
\frac16\bigr)$。这些多项式 $1$、$X - \frac12$、$X^2 - X +
\frac16$ 在相差常数倍的意义下，依次是前几个\emph{勒让德多项式}，
它们定义在区间 $\intcc{0}{1}$ 上；构造
逐次升高一个次数，每个新多项式都与
其全部前面的多项式正交。注意算法如何复用先前的工作：
第三步中减去的投影，恰好是
$X^2$ 的最佳仿射逼近，正是在\cref{ex:b1:euclid:bestapprox}中求得的
--- 格拉姆–施密特过程\emph{正是}迭代的正交投影。
\end{example}

\begin{theorem}[正交投影]\label{thm:b1:euclid:projection}
设 $F$ 为欧几里得空间 $E$ 的子空间，则
\[
E = F \oplus F^{\perp},
\]
相应的投影 $p_F$ 是到 $F$ 上的（即\emph{正交
投影}\index{正交投影}），在任一
标准正交基 $(e_1, \dots, e_k)$（取自 $F$）下，由 $p_F(x) = \sum_i
\langle x, e_i\rangle e_i$ 给出。它实现了到 $F$ 的距离：
对所有 $y \in F$，
\[
\norm{x - p_F(x)} \leq \norm{x - y},
\]
等号仅当 $y = p_F(x)$ 时成立；记 $d(x, F) = \norm{x -
p_F(x)}$。
\end{theorem}

\begin{proof}
取 $(e_i)_{i \leq k}$ 为 $F$ 的一组标准正交基
（见\cref{thm:b1:euclid:gramschmidt}，在 $F$ 中）并令 $\pi(x) =
\sum \langle x, e_i\rangle e_i \in F$。于是 $x - \pi(x) \perp e_j$
对每个 $j$ 成立（与上面的消去相同），因此 $x - \pi(x) \in
F^\perp$，即 $E = F + F^\perp$。又 $F \cap F^\perp = \{0\}$：这样的
向量满足 $\langle x, x\rangle = 0$。于是和为直和，且
$\pi = p_F$。

距离：对 $y \in F$，分解 $x - y = (x - p_F(x)) + (p_F(x) -
y)$，两项正交（分别属于 $F^\perp$ 与 $F$）；由毕达哥拉斯定理：
\[
\norm{x - y}^2 = \norm{x - p_F(x)}^2 + \norm{p_F(x) - y}^2
\geq \norm{x - p_F(x)}^2,
\]
等号成立当且仅当 $y = p_F(x)$。
\end{proof}

\begin{example}[投影不会使长度变长]
\label{ex:b1:euclid:bessel}
对分解 $x = p_F(x) + (x - p_F(x))$ 应用毕达哥拉斯定理：
\[
\norm{p_F(x)}^2 = \norm x^2 - \norm{x - p_F(x)}^2 \leq
\norm x^2 ,
\]
等号成立当且仅当 $x \in F$。取 $(e_1,
\dots, e_k)$ 为 $F$ 的标准正交基，则上式化为 $\sum_{i \leq k}\langle x,
e_i\rangle^2 \leq \norm x^2$（即\emph{贝塞尔不等式}）：
无论测量多少个标准正交方向，平方坐标
之和永远不超过平方长度 --- 对照
\cref{ex:b1:euclid:parsevalcheck}中的精确等式（当
族为完整基时）。这一行的不等式正是使
傅里叶系数在本科三年级分册中可求和的原因；此处它
已足以说明：在最小二乘拟合中加入更多基函数
只会使残差减小。
\end{example}

\begin{example}[最佳二次逼近]\label{ex:b1:euclid:bestapprox}
在 $C(\intcc{0}{1})$ 上取 $\langle f, g\rangle = \int_0^1 fg$，
则次数 $\leq 1$ 的多项式中，与 $f(x) = x^2$ 最接近的那一个，
在相应的（二次平均）距离下为 $p_F(f)$，其中 $F = \R_1[X]$。
对 $(1, X)$ 作格拉姆–施密特正交化：$e_1 = 1$，$w_2 = X - \frac12$，
$\norm{w_2}^2 = \int_0^1 (x - \frac12)^2 = \frac{1}{12}$，$e_2 =
\sqrt{12}\,(X - \tfrac12)$。于是
\[
p_F(f) = \langle f, e_1\rangle e_1 + \langle f, e_2\rangle e_2
= \frac13 + 12\Bigl(\int_0^1 x^2\bigl(x - \tfrac12\bigr)\dd
x\Bigr)\bigl(X - \tfrac12\bigr)
= X - \frac{1}{6},
\]
利用 $\int_0^1 x^2(x - \frac12)\dd x = \frac14 - \frac16 =
\frac{1}{12}$。“最小二乘”的思想，
用线性代数中的一行即可表达。
\end{example}

\begin{method}[求距离 $d(x, F)$ 的三种途径]
\label{met:b1:euclid:distance}
\begin{enumerate}
  \item \emph{$F$ 的标准正交基}：则 $p_F(x) = \sum_i
        \langle x, e_i\rangle e_i$，并由毕达哥拉斯定理，
        \[
        d(x, F)^2 = \norm{x}^2 - \norm{p_F(x)}^2
        = \norm x^2 - \sum_i \langle x, e_i\rangle^2 ,
        \]
        通常比直接计算 $x - p_F(x)$ 省事。
  \item \emph{正规方程}：取 $F$ 的任一生成族，
        解 $\langle x - p, v_j\rangle = 0$ 以求 $p$ 的
        系数（见\cref{exo:b1:euclid:5}）——无需
        标准正交化。
  \item \emph{借助补空间}：若 $F^\perp$ 的维数比 $F$ 小
        （例如 $F$ 为超平面，$F^\perp$ 为直线
        $\operatorname{Vect}(n)$），则改为
        向 $F^\perp$ 投影，即：
        \[
        d(x, F) = \norm{p_{F^\perp}(x)} =
        \frac{\abs{\langle x, n\rangle}}{\norm n} ,
        \]
        这正是经典的点到平面距离公式
        （\cref{exo:b1:multivar:8} 用到了它）。
\end{enumerate}
第 3 种途径是一般思路的特例：总是投影到
$F$、$F^\perp$ 中维数较小者。
\end{method}

\begin{example}[三角函数的正交性：傅里叶预告]
\label{ex:b1:euclid:trigortho}
在 $C(\intcc{0}{2\pi})$ 上取 $\langle f, g\rangle =
\frac1\pi\int_0^{2\pi} fg$，则函数族
\[
\Bigl(\frac{1}{\sqrt2},\ \cos x,\ \sin x,\ \cos 2x,\ \sin 2x,\
\dots\Bigr)
\]
是标准正交的：例如 $\langle \cos px, \cos qx\rangle =
\frac1\pi\int_0^{2\pi}\cos px\cos qx\,\dd x = 0$，当 $p \neq q$ 时
（将乘积化为 $\frac12[\cos(p{-}q)x +
\cos(p{+}q)x]$，并在整周期上积分），而
$\frac1\pi\int_0^{2\pi}\cos^2 px\,\dd x = 1$。
这些函数中前 $2N + 1$ 个张成的子空间上的正交投影
其坐标为 $\langle f, e_i\rangle$ ——即与余弦、正弦
的积分。这些即\emph{傅里叶
系数}，而该投影即为 $f$ 的最佳
平方平均三角逼近；本科三年级分册
研究它们的收敛性。正交性起了全部作用：
系数的公式恰是
\cref{thm:b1:euclid:projection}中的公式，原样不变。
\end{example}

\section{平面上的等距变换}

\begin{definition}\label{def:b1:euclid:isometry}
称欧几里得空间的自同态 $u$ 为
\emph{等距变换}\index{等距变换}（或称正交映射），当它保持
范数时，即 $\norm{u(x)} = \norm x$ 对所有 $x$ 成立 ——
等价地（由极化恒等式）它保持内积；
等价地，它在标准正交基下的矩阵 $A$ 满足
$A^{\mathsf T} A = I$。等距变换构成群，即正交群
$O(E)$。
\end{definition}

\begin{example}[一眼识别等距变换]
\label{ex:b1:euclid:recognize}
判断 $A = \dfrac15\begin{pmatrix} 3 & -4\\ 4 & 3\end{pmatrix}$
是否正交？列向量的范数为 $\frac15\sqrt{9 + 16} = 1$ 与
$\frac15\sqrt{16 + 9} = 1$；两列的内积为 $\frac1{25}(3\cdot(-4) +
4\cdot3) = 0$。是正交矩阵 --- 且 $\det A = \frac{9 + 16}{25} = 1$，故
它是旋转 $R_\theta$，其中 $\cos\theta = \frac35$，
$\sin\theta = \frac45$（即“$3$-$4$-$5$ 旋转”，其
角度并非 $\pi$ 的某个特殊分数）。反之，$B =
\frac{1}{\sqrt2}\begin{pmatrix} 1 & 1\\ 0 & 1\end{pmatrix}$ 看上去
像是行列式为单位的矩阵，但第一列
并非单位向量（$\frac1{\sqrt2}$）：不是正交矩阵 --- 仅凭行列式为 $\pm1$
不足以说明任何问题，必须检验各列。
\end{example}

\begin{theorem}[平面上的等距变换]\label{thm:b1:euclid:plane}
在欧几里得平面的标准正交基下，等距变换的矩阵
恰为
\[
R_\theta = \begin{pmatrix}
\cos\theta & -\sin\theta\\ \sin\theta & \cos\theta
\end{pmatrix}
\quad (\text{旋转角为 } \theta,\ \det = 1),
\]
\[
S_\theta = \begin{pmatrix}
\cos\theta & \sin\theta\\ \sin\theta & -\cos\theta
\end{pmatrix}
\quad (\det = -1),
\]
后者是关于与第一个基向量夹角为
$\frac\theta2$ 的直线的反射。
\end{theorem}

\begin{proof}
设 $A = \begin{pmatrix} a & c\\ b & d\end{pmatrix}$ 满足
$A^{\mathsf T}A = I$：各列为单位向量且相互正交。第一
列为 $(\cos\theta, \sin\theta)$，其中 $\theta$ 为某个实数；第二
列为与之正交的单位向量，即 $\pm(-\sin\theta, \cos\theta)$。符号
$+$ 时得 $R_\theta$，$-$ 时得 $S_\theta$。易验证
$S_\theta^2 = I$，且向量 $(\cos\frac\theta2,
\sin\frac\theta2)$ 不动，而与之正交的向量被取反，是
一个反射。（另有 $R_\alpha R_\beta = R_{\alpha+\beta}$：
旋转群即角度群 --- 参见
\cref{thm:b1:complex:funceq}。）
\end{proof}

\begin{omfigure}
\begin{tikzpicture}[scale=1.15]
  \draw[->, thin] (-1.2,0) -- (2.9,0) node[below] {$x$};
  \draw[->, thin] (0,-1.2) -- (0,2.6) node[left] {$y$};
  \draw[omDef, very thick] (-0.9,0) -- (2.6,0)
    node[below, pos=0.9, font=\small] {$r_1$ 的轴};
  \draw[omThm, very thick] (-0.9,-0.9) -- (2.3,2.3)
    node[right, pos=0.97, font=\small] {$r_2$ 的轴};
  \fill (2,0.5) circle (0.05) node[right, font=\small] {$M$};
  \fill (2,-0.5) circle (0.05)
    node[right, font=\small] {$r_1(M)$};
  \fill (-0.5,2) circle (0.05)
    node[above, font=\small] {$r_2(r_1(M))$};
  \draw[omDef, dashed] (2,0.5) -- (2,-0.5);
  \draw[omThm, dashed] (2,-0.5) -- (-0.5,2);
  \draw[omProp, thick, ->] (0.9,0) arc (0:45:0.9);
  \node[omProp, font=\small] at (1.15,0.32) {$\frac\pi4$};
\end{tikzpicture}

{\small 两次反射构成一个旋转：将 $M = (2,
0.5)$ 先关于 $x$ 轴反射，再关于直线 $y = x$ 反射，落在
$(-0.5, 2)$ 处 --- 即 $M$ 经旋转角为
$\frac\pi2$ 的绕原点旋转下的像，其角度恰为\emph{两倍}
于两轴夹角 $\frac\pi4$。周末问题将这一图景
推广为平面上所有等距变换的复合法则。}
\end{omfigure}

\begin{remark}[常见误区]
\label{rem:b1:euclid:pitfalls}
\emph{投影公式需要标准正交基}：若
$(v_i)$ 仅为 $F$ 的生成组，则和式 $\sum_i\langle
x, v_i\rangle v_i$ \emph{并不}是 $p_F(x)$（取 $F = \R^2$，$v_1
= e_1$，$v_2 = e_1 + e_2$）；若组非标准正交，则
改为求解正规方程
（见\cref{met:b1:euclid:distance}~(2)）。
\emph{正交组要线性无关，就不能含 $0$}：零
向量与一切向量正交，包括它自身 ---
在\cref{prop:b1:euclid:orthfree}中，线性无关要求
各向量非零。
\emph{$F^\perp$ 依赖于内积}：在 $\R_1[X]$ 中，
$\operatorname{Vect}(X)$ 的正交补，对于 $\int_0^1 PQ$ 而言，
不是常数，而是 $\operatorname{Vect}(1 - \frac32 X)$ ---
计算 $\int_0^1 x(1 - \frac32 x) = \frac12 - \frac12 = 0$；
“垂直”在内积确定之前毫无意义。
\emph{不要对 $\norm{x + y}$ 作线性展开}：正确的
恒等式是 $\norm{x+y}^2 = \norm x^2 + 2\langle x, y\rangle +
\norm y^2$；交叉项仅在正交时为零
（即毕达哥拉斯定理），而三角不等式是
\emph{不等式}。
\emph{把单位向量映为单位向量并不够}：$u(x,
y) = (x + y,\ 0)$ 把两个标准基向量都映为
单位向量 $(1, 0)$，但 $\norm{u(1,1)} = 2 \neq \sqrt2$，故
不是等距变换。定义要求 $\norm{u(x)} = \norm x$ 对
\emph{所有} $x$ 成立；用矩阵语言即 $A^{\mathsf T}A = I$，亦即
各列为单位向量\emph{且两两正交} --- 两个
条件需一并检验。
\end{remark}

\begin{remark}[内积的应用]
\label{rem:b1:euclid:whereused}
正交投影是本章应用最广的定理：
它是最小二乘法的基础（
\cref{ch:b1:multivar}中的周末问题在其上构造回归直线），傅里叶
系数（\cref{ex:b1:euclid:trigortho}），以及
\cref{exo:b1:euclid:5}中的正规方程，数值分析
大规模求解这些方程。$R_\theta / S_\theta$ 的分类将在下文
完成：周末问题将对平面上\emph{全部}
等距变换（无论是否线性）进行分类，
并研究它们的有限群 --- 花饰与
正多边形的数学。在本科二年级分册中，内积与
特征值理论相遇（对称矩阵、二次型）；在
本科三年级中，无穷维欧几里得几何化为希尔伯特
空间理论。
\end{remark}

\begin{remark}[第 3 册内的几个视角]
\label{rem:b1:euclid:perspectives}
本章延伸出两座桥。\emph{向后}，通向线性
代数：\cref{exo:b1:euclid:11}中的格拉姆矩阵把内积
打包进\cref{ch:b1:det}中的行列式理论，
而正交投影是\cref{ch:b1:linmaps}中的特殊投影算子，
其核为
$F^\perp$ --- 其全部代数（$p^2 = p$，$s = 2p -
\mathrm{id}$）原样适用，且此时还有额外收获：$\norm{x
- p(x)}$ 是一个\emph{距离}。\emph{向前}，通向分析：
\cref{ch:b1:curves}借助本章的范数度量弧长，
若缺少本章的等距变换，便无从做分类；
\cref{ch:b1:multivar}通过柯西–施瓦茨不等式读出梯度
（最速上升），并以
最小二乘法收尾本分册，这即是把\cref{thm:b1:euclid:projection}用于
$\R^n$ 中的一个数据向量。内积
正是本书的代数与分析终于相遇之处。
\end{remark}

\section{练习}

\begin{exercise}[$\star$]\label{exo:b1:euclid:1}
在 $\R^3$ 中取标准内积，计算 $\langle u, v\rangle$、$\norm u$、
$\norm v$，以及 $u = (1, 2, 2)$ 与 $v = (2, -2,
1)$ 之间的夹角。用数值方法验证柯西–施瓦茨不等式。
\end{exercise}

\begin{exercise}[$\star$]\label{exo:b1:euclid:2}
证明平行四边形恒等式 $\norm{x+y}^2 + \norm{x-y}^2 =
2\norm x^2 + 2\norm y^2$（对任意欧几里得空间），并利用它说明
$\R^2$ 上的无穷范数 $\norm{(x,y)}_\infty = \max(\abs x,
\abs y)$ 不能由任何内积导出。
\end{exercise}

\begin{exercise}[$\star$]\label{exo:b1:euclid:3}
对 $\bigl((1,1,0), (1,0,1), (0,1,1)\bigr)$，在标准
$\R^3$ 中施行格拉姆–施密特正交化。
\end{exercise}

\begin{exercise}[$\star$]\label{exo:b1:euclid:4}
在 $\R^3$ 中，设 $F = \operatorname{Vect}\bigl((1,1,1)\bigr)$。
求 $F^\perp$（方程与基）、$p_F$ 在
标准基下的矩阵，以及 $d\bigl((1, 2, 3), F\bigr)$。
\end{exercise}

\begin{exercise}[$\star\star$]\label{exo:b1:euclid:5}
（正规方程）设 $F =
\operatorname{Vect}\bigl((1,0,1),(0,1,1)\bigr) \subseteq \R^3$，且
$x = (1, 1, 4)$。求 $p_F(x)$，方法是解 $\langle x - p, v
\rangle = 0$，对两个生成元（$p = \alpha(1,0,1) +
\beta(0,1,1)$），再求 $d(x, F)$。
为何这里不需要格拉姆–施密特正交化？
\end{exercise}

\begin{exercise}[$\star\star$]\label{exo:b1:euclid:6}
设 $f$ 在 $\intcc{0}{1}$ 上连续，证明
\[
\Bigl(\int_0^1 f\Bigr)^{\!2} \leq \int_0^1 f^2 ,
\]
并讨论等号成立的情形，它是柯西–施瓦茨不等式在
$C(\intcc{0}{1})$ 中的一个实例。然后证明 $\bigl(\sum_{i=1}^n a_i\bigr)^2 \leq
n \sum a_i^2$，对实数 $a_i$ 成立。
\end{exercise}

\begin{exercise}[$\star\star$]\label{exo:b1:euclid:7}
证明：对欧几里得空间的任一子空间 $F$，
$(F^{\perp})^{\perp} = F$ 且 $\dim F^{\perp} = \dim E - \dim F$。
\end{exercise}

\begin{exercise}[$\star\star$]\label{exo:b1:euclid:8}
确定下列矩阵所对应的平面等距变换
\[
A = \frac{1}{\sqrt 2}\begin{pmatrix} 1 & -1\\ 1 & 1\end{pmatrix},
\qquad
B = \frac{1}{5}\begin{pmatrix} 3 & 4\\ 4 & -3\end{pmatrix}
\]
（指出其类型，以及旋转角或对称轴）。不做矩阵乘法，
计算 $A^8$ 与 $B^2$。
\end{exercise}

\begin{exercise}[$\star\star\star$]\label{exo:b1:euclid:9}
（将最小值视为投影）计算
\[
\min_{(a, b) \in \R^2} \int_0^1 \bigl(x^2 - a - bx\bigr)^2 \dd x ,
\]
利用\cref{ex:b1:euclid:bestapprox}：该最小值为 $\norm{f -
p_F(f)}^2$，其中 $f = X^2$，$F = \R_1[X]$。
\end{exercise}

\begin{exercise}[$\star\star\star$]\label{exo:b1:euclid:10}
设 $u$ 为欧几里得空间 $E$ 的等距变换。证明
$\ker(u - \mathrm{id}) \perp \operatorname{im}(u - \mathrm{id})$，
并由此推出 $E = \ker(u - \mathrm{id}) \oplus \operatorname{im}(u -
\mathrm{id})$。\emph{（计算 $\langle x - u(x), y\rangle$，
其中 $u(y) = y$，并利用内积保持不变的性质。）}
\end{exercise}

\begin{exercise}[$\star\star$]\label{exo:b1:euclid:11}
（格拉姆矩阵）对欧几里得空间中的向量 $v_1, \dots, v_k$，
记 $G = \bigl(\langle v_i, v_j\rangle\bigr)_{1 \leq i, j
\leq k}$ 为它们的\emph{格拉姆矩阵}。
\begin{enumerate}
  \item 证明 $(v_1, \dots, v_k)$ 线性无关当且仅当
        $G$ 可逆。\emph{（若 $Gc = 0$，计算
        $\norm{\sum_i c_i v_i}^2$。）}
  \item 计算 $(1, X, X^2)$ 在 $\R_2[X]$ 中的格拉姆矩阵，
        取 $\langle P, Q\rangle = \int_0^1 PQ$，并认出
        希尔伯特矩阵 $H_3$，它出现于
        \cref{ch:b1:det}的周末问题中，并由 $\det H_3 =
        \frac1{2160} \neq 0$ 得出线性无关。
\end{enumerate}
\end{exercise}

\begin{exercise}[$\star\star\star$]\label{exo:b1:euclid:12}
设 $u$ 为欧几里得空间 $E$ 上的自同态，它在某组标准正交基下的矩阵
$A$ 既是正交矩阵（$A^{\mathsf T}A
= I$），又是对称矩阵（$A^{\mathsf T} = A$）。
\begin{enumerate}
  \item 证明 $A^2 = I$，并借助
        \cref{thm:b1:linmaps:projchar}推出 $E = \ker(u -
        \mathrm{id}) \oplus \ker(u + \mathrm{id})$。
  \item 证明两个子空间正交，从而 $u$
        是\emph{正交对称}，关于 $F =
        \ker(u - \mathrm{id})$：即关于 $F$ 的对称。
        \emph{（对 $u(x) = x$ 与 $u(y) = -y$，两种方式计算 $\langle
        x, y\rangle$。）}
  \item 对平面情形分类：在
        \cref{thm:b1:euclid:plane}中的哪些矩阵是对称的，
        以及它们对应的映射是什么？
\end{enumerate}
\end{exercise}

\section{问题：平面等距变换与莱昂纳多定理}

\begin{problem}\label{pb:b1:euclid:1}
\emph{平面等距变换}是指任一映射 $f \colon \R^2 \to
\R^2$ 保持距离，即 $\norm{f(x) - f(y)} = \norm{x - y}$
对所有 $x, y$ 成立 --- 不假设线性。本题证明
这些映射恰为平移、旋转、
对称与滑动对称（平面等距变换的
分类），计算它们的复合，并确定
全部\emph{有限}群：莱昂纳多定理，
即花饰图案背后的数学。自第 II 部分起，将 $\R^2$ 与
$\C$ 等同（见\cref{ch:b1:complex}）：标准内积为
$\langle z, w\rangle = \operatorname{Re}(z\conj w)$，而范数则
为模。\index{等距变换}\index{滑动对称}
\index{莱昂纳多定理}\index{二面体群}

\medskip
\noindent\textbf{第 I 部分 --- 每个等距变换都是仿射的。}
\begin{enumerate}
  \item 验证平移 $t_a(x) = x + a$、线性等距变换，
        以及它们的全部复合都是等距变换，
        并且等距变换在复合运算下构成群。
  \item 设 $f$ 为等距变换，且 $f(0) = 0$。证明 $f$
        保持范数，再 --- 借助极化，
        由\cref{thm:b1:euclid:cs} --- 可得 $\langle f(x),
        f(y)\rangle = \langle x, y\rangle$，对所有 $x, y$ 成立。
  \item 仍设 $f(0) = 0$，展开 $\norm{f(x + y) - f(x) -
        f(y)}^2$ 与 $\norm{f(\lambda x) - \lambda f(x)}^2$，
        利用第 2 小问，断定 $f$ 是
        \emph{线性}的：$f \in O(\R^2)$。
  \item 由此推出每个等距变换 $f$ 都可\emph{唯一}地写成
        $f = t_a \circ g$，其中 $a = f(0)$，$g$ 是
        线性等距变换（$f$ 的\emph{线性部分}）。
  \item 称 $f$ 为\emph{保向}的（若 $\det g = 1$），\emph{反向}
        的（若 $\det g = -1$）。证明复合映射的线性部分
        是各线性部分的复合，
        并陈述由此得到的符号法则
        （保向与反向的复合规律与 $+1/-1$ 的乘法相同）。
\end{enumerate}

\medskip
\noindent\textbf{第 II 部分 --- 四种类型。}借助
\cref{thm:b1:euclid:plane}，$\C$ 上的线性等距变换为 $z
\mapsto az$ 与 $z \mapsto a\conj z$，其中 $\abs a = 1$；于是每个平面等距变换
可写成
\[
f(z) = a z + b \quad (\text{保向}) \qquad\text{或}\qquad
f(z) = a\conj z + b \quad (\text{反向}), \qquad \abs a = 1 .
\]
\begin{enumerate}[resume]
  \item 验证对应关系：$R_\theta$ 为 $z \mapsto
        \eu^{\iu\theta}z$，$S_\theta$ 为 $z \mapsto
        \eu^{\iu\theta}\conj z$（分别对 $1$ 与 $\iu$ 验证）。
  \item （保向情形）设 $f(z) = az + b$，$\abs a = 1$。证明：
        若 $a = 1$，则 $f$ 为平移；若 $a \neq 1$，则 $f$ 有
        唯一不动点 $z_0 = b/(1 - a)$，且 $f(z) - z_0
        = a(z - z_0)$：即中心为 $z_0$ 的\emph{旋转}，
        旋转角为 $\arg a$。
  \item （反向情形）设 $f(z) = a\conj z + b$，$v =
        a\conj b + b$。证明 $f \circ f = t_v$，且 $f \circ t_v =
        t_v \circ f$。若 $v = 0$，证明 $z$ 与
        $f(z)$ 的中点为不动点，并证明 $f$ 是
        关于某直线的\emph{对称}。若 $v \neq 0$，证明 $r
        = t_{-v/2}\circ f$ 是一个对称，其轴
        平行于 $v$，从而 $f = t_{v/2} \circ r$ 是一个
        \emph{滑动对称}。综上，每个平面等距变换
        都是平移、旋转、对称或滑动对称
        （平面等距变换的分类）。
  \item （复合）证明：转角分别为
        $\alpha$ 与 $\beta$ 的两个旋转的复合是转角为
        $\alpha + \beta$ 的旋转（当 $\alpha + \beta \in
        2\pi\Z$ 时为平移）；两个对称的复合是
        转角为两轴夹角两倍的旋转（
        两轴平行时为平移）。
  \item 由此推出每个平面等距变换都是至多
        三个对称的复合。
\end{enumerate}

\medskip
\noindent\textbf{第 III 部分 --- 三项辨识。}
\begin{enumerate}[resume]
  \item 完全分类 $f(z) = \iu\conj z + 1 + \iu$：
        类型、轴与滑动向量。
  \item 设 $f$ 为转角为 $\frac\pi2$、中心为 $0$ 的旋转，
        而 $g$ 为转角为 $\frac\pi2$、中心为 $1$ 的旋转。
        求 $g \circ f$ 的形式 $z \mapsto az + b$，并
        判定其类型、中心与转角。
  \item 设 $r_1(z) = \conj z$（关于实轴的对称），而
        $r_2(z) = \iu\conj z$（关于直线 $y = x$ 的对称）。
        求 $r_2 \circ r_1$，并在此例中验证第 9 小问
        的结论。
\end{enumerate}

\medskip
\noindent\textbf{第四部分 --- 有限群：达·芬奇定理。}
设 $G$ 为平面等距变换的\emph{有限}群。
\begin{enumerate}[resume]
  \item 证明等距变换保持重心：若
        $\lambda_1 + \dots + \lambda_m = 1$ 且 $f = t_a \circ
        g$（$g$ 为线性映射），则 $f\bigl(\sum_i \lambda_i
        x_i\bigr) = \sum_i \lambda_i f(x_i)$。
  \item 令 $c = \frac{1}{\abs G}\sum_{g \in G} g(x_0)$，
        其中 $x_0$ 为任取的点。证明每个 $h \in G$ 都固定 $c$：
        即有限等距变换群有公共不动点。
  \item 由此可知，经 $t_{-c}$ 共轭后，
        可假设 $G \subseteq O(\R^2)$。令 $G^{+} = \{g \in G :
        \det g = 1\}$。证明或者 $G = G^{+}$，或者 $G^{+}$
        的指数恰为 $2$（在 $G$ 中；给出一个双射 $G^+
        \to G \setminus G^+$）。
  \item 证明绕 $c$ 的\emph{旋转}构成的有限群
        是循环群：在其中取转角最小的旋转
        $R_{\theta_0}$，其转角 $\theta_0 \in
        \intoo{0}{2\pi}$。利用角的带余除法
        证明它生成该群。于是得 $\theta_0 =
        \frac{2\pi}n$，且 $G^{+} = \{R_{\theta_0}^{\,k}\} \cong
        C_n$。
  \item 设 $G \neq G^{+}$，并取一个反射 $s \in G$。
        证明 $G = G^{+} \cup sG^{+}$，且 $s r s = r^{-1}$
        对每个旋转 $r \in G^{+}$ 成立，并证明所有 $n$ 个
        属于 $sG^{+}$ 的元素均为反射。于是 $G$ 为\emph{二面体
        群} $D_n$，阶为 $2n$。
  \item 综上得（\emph{达·芬奇定理}）：平面上的每个有限等距变换群
        都是循环群 $C_n$ 或二面体群
        $D_n$。
\end{enumerate}

\medskip
\noindent\textbf{第五部分 --- 附加结论与综合。}
\begin{enumerate}[resume]
  \item 直接证明：有限等距变换群中不含
        非恒等的平移，也不含滑动反射
        （可考虑此类元素的
        幂）。
  \item 设 $P_n$ 为以 $n$ 次单位根为顶点的正 $n$ 边形
        （$n \geq 3$）。证明其对称群恰为 $D_n$：
        即 $n$ 个旋转 $z
        \mapsto \omega^k z$ 与 $n$ 个反射 $z \mapsto
        \omega^k \conj z$，其中 $\omega = \eu^{2\iu\pi/n}$，再无其他
        元素。
  \item 给出平面图形，其对称群依次为 $C_1$、
        $D_1$、$D_2$ 与 $C_3$。
  \item 列出顶点为 $\pm1, \pm\iu$ 的正方形对称群的八个元素，
        写成 $z \mapsto
        \omega^k z$ 或 $z \mapsto \omega^k\conj z$ 的形式，并给出
        四个反射中每一个的对称轴。
  \item 设 $f, g$ 为同一转角 $\theta
        \notin 2\pi\Z$ 的旋转，分别绕\emph{不同}的中心 $c_1 \neq
        c_2$。逐点计算 $f\circ g - g\circ f$，并证明
        $f\circ g \neq g\circ f$；进一步证明 $(f \circ
        g)\circ(g\circ f)^{-1}$ 为非恒等的平移，从而
        任何包含 $f$ 与 $g$ 的群都是无限群 ——
        这是对达·芬奇定理中唯一中心的第二种
        解释。
  \item 综合题，分四句作答：哪两个结构性结果可把任意等距变换
        化归为线性代数（第 3--4 题），
        以及把任意有限群化归为
        $O(2)$ 的子群（第 15 题）；平面等距变换的完整列表是什么，
        哪些不变量（保向或反向、是否有不动点）
        可区分这四种类型；为何第 9 题中的复合规则使反射
        成为一切变换的生成元；
        达·芬奇定理在有限尺度上补充了什么；
        并指出
        第二、四部分中证明的两个定理。
\end{enumerate}
\end{problem}
